\documentclass[11pt]{article}
\usepackage[margin=1in]{geometry}
\usepackage{amsmath,amssymb,mathtools,amsthm}
\usepackage{booktabs,longtable}
\usepackage{enumitem}
\usepackage[hidelinks]{hyperref}

\newcommand{\dd}{\mathrm{d}}
\newcommand{\R}{\mathbb{R}}
\newcommand{\TT}{\mathrm{TT}}
\newcommand{\E}{\mathbb{E}}
\newcommand{\calX}{\mathcal{X}}
\newcommand{\calR}{\mathcal{R}}
\newcommand{\calB}{\mathcal{B}}
\newcommand{\calM}{\mathcal{M}}
\newcommand{\calF}{\mathcal{F}}
\newcommand{\dotp}{\partial_i}

\newtheorem{proposition}{Proposition}
\newtheorem{corollary}{Corollary}
\theoremstyle{definition}
\newtheorem{definition}{Definition}
\newtheorem{assumption}{Assumption}
\theoremstyle{remark}
\newtheorem{remark}{Remark}

\title{Self-Contained Fixed-Branch Proofs for Zhao--Cui Tensor-Train Sequential Filtering}
\author{BayesFilter P12 Proof Expansion Note}
\date{2026-05-31}

\begin{document}
\maketitle

\begin{abstract}
This note expands the fixed-branch Zhao--Cui tensor-train derivative contract
into a self-contained proof document.  It first derives the ordinary Bayesian
filtering recursion, then explains the Zhao--Cui squared tensor-train
construction in BayesFilter notation, and finally introduces a deterministic
fixed-branch variant.  Two propositions are proved.  Proposition 1 proves that
the fixed-branch squared-TT recursion defines a normalized nonnegative
approximate Bayesian filter.  Proposition 2 proves that the analytical gradient
is the exact derivative of the declared approximate scalar
\[
  \widehat\ell_T^{\TT}(\alpha)
  =
  \sum_{t=1}^T\{\log \widehat z_t(\alpha)-c_t(\alpha)\}.
\]
The proofs are for the declared fixed branch only.  They do not prove posterior
accuracy, HMC convergence, or differentiability of the adaptive TT-cross,
rank-changing, sample-changing companion code.
\end{abstract}

\section{Introduction And Reader Contract}

The problem is sequential inference in a state-space model.  At each time
\(t\), one wants a filtering density over the current state and unknown
parameters after seeing data \(y_{1:t}\).  The exact density is usually not
available in closed form.  Zhao and Cui propose to approximate the
high-dimensional density that appears inside the filtering recursion by a
tensor-train (TT) representation of a square-root density, then square that TT
and add a small defensive term so the resulting approximation is nonnegative.
The squared representation also gives marginal densities and conditional
Knothe--Rosenblatt (KR) maps.

This note separates three statements.

\begin{enumerate}[leftmargin=*]
\item The Zhao--Cui construction is a source-grounded high-dimensional
sequential approximation strategy.  This is supported by their Algorithm 1,
their squared-TT construction around their equation (13), Algorithm 2, and
their discussion of the evidence normalizer in Section 4.1
\cite{zhao-cui-2024}.
\item The BayesFilter fixed-branch variant is a mathematically declared
approximate filtering recursion.  The proof is a project derivation: it shows
that the output densities are nonnegative and normalized.
\item The fixed-branch analytical gradient differentiates the same approximate
scalar that the fixed-branch recursion computes.  This is also a project
derivation.  It requires fixed TT ranks, fixed interpolation or least-squares
systems, fixed basis maps, fixed random/debug arrays, fixed rounding choices,
fixed QR/SVD/Cholesky branches, and a declared rule for the stabilizing
constant.
\end{enumerate}

The note does not prove that the approximation equals the exact posterior
filter.  It proves correctness of the approximate recursion and correctness of
the gradient for the declared approximate scalar.  Accuracy must be studied by
approximation theory or numerical experiments.

\section{State-Space Filtering From First Principles}

Let \(x_t\in\calX_t\subseteq\R^{d_x}\) be the state at time \(t\), let
\(\vartheta\in\Theta\subseteq\R^{d_\vartheta}\) be a static parameter, and let
\(y_t\) be the observed datum.  A dominated state-space model specifies

\[
  p_0(x_0,\vartheta),\qquad
  f_\alpha(x_t\mid x_{t-1},\vartheta),\qquad
  g_\alpha(y_t\mid x_t,\vartheta),
\]

where \(\alpha\) denotes an external model parameter with respect to which a
future likelihood-gradient calculation may differentiate.  The parameter
\(\vartheta\) is a random coordinate in the filtering distribution; \(\alpha\)
is the coordinate of the scalar derivative in this note.

Assume that the filter at time \(t-1\) has density
\[
  p_{t-1}(x_{t-1},\vartheta\mid y_{1:t-1})
\]
with respect to \(\dd x_{t-1}\dd\vartheta\).  The prediction density is obtained
by integrating over the previous state:
\[
  p(x_t,\vartheta\mid y_{1:t-1})
  =
  \int
  f_\alpha(x_t\mid x_{t-1},\vartheta)
  p_{t-1}(x_{t-1},\vartheta\mid y_{1:t-1})
  \,\dd x_{t-1}.
  \label{eq:p12-prediction}
\]
Bayes' rule then gives the unnormalized update
\[
  \tilde p_t(x_t,\vartheta)
  =
  g_\alpha(y_t\mid x_t,\vartheta)
  p(x_t,\vartheta\mid y_{1:t-1}).
\]
Its integral
\[
  Z_t(\alpha)
  =
  \int
  g_\alpha(y_t\mid x_t,\vartheta)
  f_\alpha(x_t\mid x_{t-1},\vartheta)
  p_{t-1}(x_{t-1},\vartheta\mid y_{1:t-1})
  \,\dd x_{t-1}\dd x_t\dd\vartheta
  \label{eq:p12-exact-evidence-increment}
\]
is the one-step evidence increment \(p_\alpha(y_t\mid y_{1:t-1})\).  The exact
filter is therefore
\[
  p_t(x_t,\vartheta\mid y_{1:t})
  =
  \frac{
  g_\alpha(y_t\mid x_t,\vartheta)
  \int f_\alpha(x_t\mid x_{t-1},\vartheta)
  p_{t-1}(x_{t-1},\vartheta\mid y_{1:t-1})\,\dd x_{t-1}
  }
  {Z_t(\alpha)}.
  \label{eq:p12-exact-filter}
\]

It is often useful to keep the previous state inside the object being
approximated.  Define
\[
  u_t=(x_t,\vartheta,x_{t-1})
\]
and
\[
  q_t(u_t;\alpha)
  =
  p_{t-1}(x_{t-1},\vartheta\mid y_{1:t-1})
  f_\alpha(x_t\mid x_{t-1},\vartheta)
  g_\alpha(y_t\mid x_t,\vartheta).
  \label{eq:p12-exact-joint-q}
\]
Then \(Z_t(\alpha)=\int q_t(u_t;\alpha)\,\dd u_t\), and the exact filter is
the marginal of the normalized joint density \(q_t/Z_t\) over the old state:
\[
  p_t(x_t,\vartheta\mid y_{1:t})
  =
  \int
  \frac{q_t(x_t,\vartheta,x_{t-1};\alpha)}{Z_t(\alpha)}
  \,\dd x_{t-1}.
  \label{eq:p12-exact-filter-marginal}
\]
This is the mathematical reason Zhao--Cui approximate the nonseparable density
over \((x_t,\vartheta,x_{t-1})\): once this joint object can be represented and
integrated, filtering and evidence increments follow by normalization and
marginalization.

\section{Tensor Trains From First Principles}

A tensor train is a separated representation of a high-dimensional function.
Let \(r=(r_1,\ldots,r_D)\) and let \(\{b_{k,a}\}_{a=1}^{n_k}\) be basis
functions for the \(k\)-th coordinate.  A scalar function \(\phi(r)\) has a
functional TT representation if
\[
  \phi(r)
  =
  G_1(r_1)G_2(r_2)\cdots G_D(r_D),
  \label{eq:p12-tt-evaluation}
\]
where
\[
  G_k(r_k)\in\R^{\rho_{k-1}\times\rho_k},\qquad
  \rho_0=\rho_D=1.
\]
The integers \(\rho_k\) are TT ranks.  Each matrix-valued core can be expanded
in the one-dimensional basis:
\[
  G_k(r_k)
  =
  \sum_{a=1}^{n_k} C_{k,a} b_{k,a}(r_k),
  \qquad
  C_{k,a}\in\R^{\rho_{k-1}\times\rho_k}.
  \label{eq:p12-core-expansion}
\]
Evaluation of \(\phi(r)\) is a chain of matrix multiplications, rather than a
full tensor-product sum over \(n_1\cdots n_D\) coefficients.

If \(\nu=\nu_1\otimes\cdots\otimes\nu_D\) is the reference measure used by the
basis, define one-dimensional mass matrices
\[
  M_k[a,b]
  =
  \int b_{k,a}(r_k)b_{k,b}(r_k)\,\dd\nu_k(r_k).
  \label{eq:p12-mass-matrix}
\]
The integral of \(\phi^2\) is obtained by the backward contractions
\[
  R_D=1,\qquad
  R_{k-1}
  =
  \sum_{a,b=1}^{n_k}
  M_k[a,b]\,C_{k,a}R_kC_{k,b}^\top,
  \qquad k=D,\ldots,1.
  \label{eq:p12-mass-contraction}
\]
Then \(R_0=\int \phi(r)^2\,\dd\nu(r)\).  The contraction is exact for the
declared TT cores and mass matrices; approximation enters through the
construction of the cores, not through this algebraic contraction.

\section{Zhao--Cui Sequential Squared-TT Algorithm}

Zhao--Cui's Algorithm 1 constructs, at each time \(t\), a nonseparable density
over \((x_t,\vartheta,x_{t-1})\), approximates it in TT form, integrates the
approximation to get a normalizer, and marginalizes to obtain the next
filtering object \cite{zhao-cui-2024}.  In their notation, the core density is
the analogue of
\[
  q_t(x_t,\vartheta,x_{t-1})
  =
  \widehat\pi_{t-1}(x_{t-1},\vartheta\mid y_{1:t-1})
  f(x_t\mid x_{t-1},\vartheta)
  g(y_t\mid x_t,\vartheta).
  \label{eq:p12-zc-q}
\]
The exact integral of \(q_t\) is the evidence increment for the exact filter if
\(\widehat\pi_{t-1}\) is exact; with \(\widehat\pi_{t-1}\) approximate, it is
the evidence increment for the declared approximate recursion.

The squared-TT construction approximates the square root of an unnormalized
nonnegative density by a TT:
\[
  \sqrt{q_t(u_t)}
  \approx
  \phi_t(u_t)
  =
  G_{t,1}(u_{t,1})\cdots G_{t,D}(u_{t,D}).
\]
Instead of using a sign-changing TT density directly, Zhao--Cui square the TT
and add a defensive density:
\[
  \widehat q_t(u_t)
  =
  \phi_t(u_t)^2+\tau_t\lambda_t(u_t),
  \qquad
  \tau_t\ge0,\quad
  \lambda_t\ge0,\quad
  \int\lambda_t=1.
  \label{eq:p12-squared-tt-density}
\]
The approximate normalizer is
\[
  \widehat Z_t
  =
  \int \widehat q_t(u_t)\,\dd u_t.
  \label{eq:p12-squared-tt-normalizer}
\]
Zhao--Cui's Section 4.1 identifies the corresponding exact joint normalizer as
evidence; their Algorithm 2 applies the squared-TT construction sequentially.
The companion code follows the same scalar path:
\texttt{models/full\_sol.m} updates
\texttt{logmarginal\_likelihood} by \texttt{log(sirt.z)-const}, while
\texttt{@TTSIRT/marginalise.m} sets
\texttt{obj.z = obj.fun\_z + obj.tau}.

The conditional KR transport is a sampling and conditioning mechanism built
from marginal and conditional densities obtained by integrating the squared TT
\cite{zhao-cui-2024,cui-dolgov-2022}.  For the two propositions below, the
transport role is secondary: we need only the nonnegative squared-TT density,
its normalizer, and the marginalization that produces the next filter.  The KR
maps explain how the same approximation can also generate samples, smoothing
paths, and conditional parameter/state draws.

\paragraph{Published-style recursion.}
At a high level, the Zhao--Cui-style squared-TT filtering recursion is:
\begin{enumerate}[leftmargin=*]
\item Start from a prior or previous approximate filter
  \(\widehat p_{t-1}(x_{t-1},\vartheta)\).
\item Form the unnormalized joint density \(q_t\) in
  \eqref{eq:p12-zc-q}.
\item Build a TT approximation \(\phi_t\approx\sqrt{q_t}\) using the selected
  basis, variable ordering, TT-cross/ALS/rank policy, and coordinate maps.
\item Define \(\widehat q_t=\phi_t^2+\tau_t\lambda_t\) and compute
  \(\widehat Z_t=\int\widehat q_t\).
\item Normalize and marginalize:
  \[
    \widehat p_t(x_t,\vartheta)
    =
    \int
    \widehat q_t(x_t,\vartheta,x_{t-1})/\widehat Z_t
    \,\dd x_{t-1}.
  \]
\item Use the marginal and conditional densities implied by the squared TT to
  build conditional KR maps when sampling, smoothing, or parameter/state draws
  are needed.
\end{enumerate}

The published and companion-code algorithms adapt many choices: cross or
interpolation points, ranks, truncations, random/debug samples, stopping
decisions, preconditioner branches, and QR/SVD signs or active subspaces.  Such
adaptation is natural for filtering.  It is not a single globally smooth map
from model parameter to likelihood scalar.  That is why the next section
declares a fixed branch.

\section{Fixed-Branch TT Filtering Variant}

\begin{definition}[Fixed branch]
A fixed branch consists of fixed state ordering, basis families, reference
domains, interpolation or cross sets, TT ranks, local solve schedule, ALS sweep
count, random/debug sample arrays, truncation decisions, QR/SVD/Cholesky signs
and active ranks, defensive density \(\lambda_t\), defensive mass rule
\(\tau_t\), coordinate maps, and the stabilizing constant rule \(c_t\).  Along
this branch, each of these choices is treated as a declared smooth object or as
a fixed numerical object.
\end{definition}

The stabilizing constant appears because the code approximates a shifted
potential.  If \(V_{t,0}\) is an unshifted potential and
\[
  V_t(r;\alpha)=V_{t,0}(r;\alpha)-c_t(\alpha),
\]
then the shifted normalizer \(\widehat z_t\) corresponds to the original-scale
increment \(\exp(-c_t)\widehat z_t\).  Thus the declared scalar is
\[
  \widehat\ell_T^{\TT}(\alpha)
  =
  \sum_{t=1}^T\{\log\widehat z_t(\alpha)-c_t(\alpha)\}.
  \label{eq:p12-declared-scalar}
\]
If a future implementation freezes \(c_t\) numerically rather than
differentiating the selected branch rule, it must name a different scalar.

\paragraph{Fixed-branch recursion.}
Given \(\widehat p_{t-1}^{\TT}\), the fixed-branch step is:
\begin{enumerate}[leftmargin=*]
\item form \(q_t\) from \(\widehat p_{t-1}^{\TT}\), \(f_\alpha\), and
  \(g_\alpha\);
\item apply the fixed coordinate/basis maps and stabilizing constant;
\item solve the fixed TT-core construction problem, giving cores
  \(G_{t,1},\ldots,G_{t,D}\);
\item define \(\widehat q_t=\phi_t^2+\tau_t\lambda_t\) and
  \(\widehat z_t=\int\widehat q_t\);
\item normalize and marginalize to obtain
  \(\widehat p_t^{\TT}\);
\item store every branch object needed to reproduce the same scalar and its
  derivative.
\end{enumerate}
The fixed branch is less adaptive than the research code.  Its purpose is not
to improve filtering accuracy.  Its purpose is to define a deterministic
approximate scalar and a same-scalar analytical gradient.

\section{Proposition 1: Fixed-Branch TT Recursion Is A Normalized Approximate Filter}

\begin{assumption}[Dominated fixed-branch recursion]
For every \(t\), all state and parameter densities are dominated by fixed
measures.  The previous approximate filter
\(\widehat p_{t-1}^{\TT}\) is nonnegative and integrates to one.  The transition
and observation densities \(f_\alpha\) and \(g_\alpha\) are nonnegative and
measurable.  The fixed branch produces a measurable square-root approximation
\(\phi_t\).  The defensive density satisfies \(\lambda_t\ge0\) and
\(\int\lambda_t=1\).  The defensive mass satisfies \(\tau_t\ge0\).  Finally,
\[
  0<\widehat z_t(\alpha)
  =
  \int\{\phi_t(u_t;\alpha)^2+\tau_t(\alpha)\lambda_t(u_t)\}\,\dd u_t
  <\infty.
\]
\end{assumption}

\begin{proposition}[Fixed-branch squared-TT recursion defines a normalized approximate filter]
Under the dominated fixed-branch recursion assumptions, define
\[
  q_t(u_t;\alpha)
  =
  \widehat p_{t-1}^{\TT}(x_{t-1},\vartheta;\alpha)
  f_\alpha(x_t\mid x_{t-1},\vartheta)
  g_\alpha(y_t\mid x_t,\vartheta),
\]
\[
  \widehat q_t(u_t;\alpha)
  =
  \phi_t(u_t;\alpha)^2+\tau_t(\alpha)\lambda_t(u_t),
  \qquad
  \widehat\pi_t(u_t;\alpha)
  =
  \frac{\widehat q_t(u_t;\alpha)}{\widehat z_t(\alpha)},
\]
and
\[
  \widehat p_t^{\TT}(x_t,\vartheta;\alpha)
  =
  \int
  \widehat\pi_t(x_t,\vartheta,x_{t-1};\alpha)
  \,\dd x_{t-1}.
  \label{eq:p12-fixed-filter-marginal}
\]
Then \(\widehat p_t^{\TT}\) is a nonnegative normalized density.  By induction
from a normalized prior, the recursion defines a sequence of normalized
approximate filtering densities.
\end{proposition}

\begin{proof}
The function \(\phi_t^2\) is nonnegative pointwise, and
\(\tau_t\lambda_t\) is nonnegative because \(\tau_t\ge0\) and
\(\lambda_t\ge0\).  Hence \(\widehat q_t\ge0\).  By assumption,
\(\widehat z_t=\int\widehat q_t\) is finite and strictly positive, so
\(\widehat\pi_t=\widehat q_t/\widehat z_t\) is well defined and nonnegative.
Moreover,
\[
  \int \widehat\pi_t(u_t;\alpha)\,\dd u_t
  =
  \frac{1}{\widehat z_t(\alpha)}
  \int \widehat q_t(u_t;\alpha)\,\dd u_t
  =
  1.
\]
Now integrate the marginal in \eqref{eq:p12-fixed-filter-marginal}.  Tonelli's
theorem applies to the nonnegative integrand, so
\[
\begin{aligned}
  \int \widehat p_t^{\TT}(x_t,\vartheta;\alpha)\,\dd x_t\dd\vartheta
  &=
  \int\int
  \widehat\pi_t(x_t,\vartheta,x_{t-1};\alpha)
  \,\dd x_{t-1}\dd x_t\dd\vartheta\\
  &=
  \int\widehat\pi_t(u_t;\alpha)\,\dd u_t
  =
  1.
\end{aligned}
\]
Thus the marginal is nonnegative and normalized.  The base case is the
normalized prior.  If \(\widehat p_{t-1}^{\TT}\) is normalized, the argument
above proves that \(\widehat p_t^{\TT}\) is normalized.  Induction over \(t\)
therefore proves the claim.
\end{proof}

\begin{corollary}[Exactness only under zero approximation error]
If, at every time \(t\), the fixed-branch squared-TT object satisfies
\(\widehat q_t(u_t;\alpha)=q_t(u_t;\alpha)\) almost everywhere and no shifted
constant changes the scale, then \(\widehat z_t=Z_t\) and the recursion
recovers the exact Bayesian filtering recursion.  Otherwise it is a normalized
approximate filtering recursion.
\end{corollary}

\begin{proof}
If \(\widehat q_t=q_t\), then their integrals agree, so
\(\widehat z_t=Z_t\).  The normalized joint density
\(\widehat q_t/\widehat z_t\) equals \(q_t/Z_t\), and marginalizing out
\(x_{t-1}\) gives the exact filter by \eqref{eq:p12-exact-filter-marginal}.
If \(\widehat q_t\ne q_t\), Proposition 1 still proves normalization and
nonnegativity, but it does not imply equality with the exact posterior.
\end{proof}

\section{Proposition 2: Same-Scalar Analytical Gradient}

\begin{assumption}[Fixed-branch differentiability]
The assumptions of Proposition 1 hold.  In addition, all branch choices are
fixed.  Along that branch, the model densities, previous approximate filter,
coordinate maps, basis maps, core coefficients, defensive mass
\(\tau_t\), and stabilizing constant \(c_t\) are differentiable in the external
parameter coordinate \(\alpha_i\).  Differentiation may be interchanged with
the finite contractions and with the displayed integrals.  Fixed interpolation
or least-squares systems used to construct TT cores are nonsingular, or a
specific differentiable pseudoinverse branch is declared.  Each normalizer
\(\widehat z_t\) is positive.
\end{assumption}

\begin{proposition}[Analytical gradient differentiates the declared fixed-branch approximate likelihood]
Under the fixed-branch differentiability assumptions, the derivative of
\[
  \widehat\ell_T^{\TT}(\alpha)
  =
  \sum_{t=1}^T
  \{\log\widehat z_t(\alpha)-c_t(\alpha)\}
\]
is
\[
  \partial_i\widehat\ell_T^{\TT}(\alpha)
  =
  \sum_{t=1}^T
  \left[
  \frac{\partial_i\widehat z_t(\alpha)}{\widehat z_t(\alpha)}
  -
  \partial_i c_t(\alpha)
  \right],
  \label{eq:p12-score-general}
\]
where
\[
  \partial_i\widehat z_t
  =
  2\int\phi_t(u;\alpha)\partial_i\phi_t(u;\alpha)\,\dd u
  +
  \partial_i\tau_t
  \label{eq:p12-z-derivative-integral}
\]
when the defensive density \(\lambda_t\) is fixed and normalized.  If
\(\lambda_t\) also depends on \(\alpha_i\), then the extra term
\(\tau_t\int\partial_i\lambda_t\,\dd u\) appears; it vanishes for normalized
fixed-support \(\lambda_t\) because \(\partial_i\int\lambda_t=0\).
Moreover, when the TT cores are produced by the displayed fixed interpolation
or fixed weighted least-squares equations, the core sensitivities are obtained
by differentiating exactly those equations, and the recursion includes the
previous-filter derivative in \eqref{eq:p12-prev-filter-derivative}.
\end{proposition}

\begin{proof}
The scalar is a finite sum of differentiable terms on the fixed branch.  Since
\(\widehat z_t>0\),
\[
  \partial_i\{\log\widehat z_t-c_t\}
  =
  \frac{\partial_i\widehat z_t}{\widehat z_t}
  -
  \partial_i c_t.
\]
Summing over \(t\) proves \eqref{eq:p12-score-general}.  Next,
\[
  \widehat z_t
  =
  \int \{\phi_t(u;\alpha)^2+\tau_t(\alpha)\lambda_t(u)\}\,\dd u.
\]
By the assumed interchange of differentiation and integration,
\[
  \partial_i\widehat z_t
  =
  \int 2\phi_t(u;\alpha)\partial_i\phi_t(u;\alpha)\,\dd u
  +
  \partial_i\tau_t\int\lambda_t(u)\,\dd u.
\]
Because \(\int\lambda_t=1\), this gives
\eqref{eq:p12-z-derivative-integral}.

It remains to derive \(\partial_i\phi_t\) from the same computation path.  If
\[
  \phi_t(u)=G_{t,1}(u_1)\cdots G_{t,D}(u_D),
\]
then repeated use of the product rule gives
\[
  \partial_i\phi_t(u)
  =
  \sum_{k=1}^D
  G_{t,1}(u_1)\cdots
  \partial_iG_{t,k}(u_k)\cdots
  G_{t,D}(u_D).
  \label{eq:p12-tt-product-rule}
\]
Write each core in a fixed basis:
\[
  G_{t,k}(u_k)
  =
  \sum_{a=1}^{n_k}C_{t,k,a}b_{k,a}(u_k).
\]
When the basis is fixed, the derivative of the core is obtained by replacing
\(C_{t,k,a}\) with \(\partial_iC_{t,k,a}\).  This is not an
extra approximation: it is the coordinate derivative of the same core
expansion.

Using the mass matrices \(M_k\) in \eqref{eq:p12-mass-matrix}, define
\[
  R_{t,D}=1,\qquad
  R_{t,k-1}
  =
  \sum_{a,b=1}^{n_k}
  M_k[a,b]\,C_{t,k,a}R_{t,k}C_{t,k,b}^\top.
  \label{eq:p12-proof-contraction}
\]
Then the square part of the normalizer is \(R_{t,0}\), and the differentiated
contraction is
\[
\begin{aligned}
  \partial_i R_{t,k-1}
  =
  \sum_{a,b=1}^{n_k}M_k[a,b]\Big[
  &(\partial_iC_{t,k,a})R_{t,k}C_{t,k,b}^\top
  +C_{t,k,a}(\partial_iR_{t,k})C_{t,k,b}^\top\\
  &+C_{t,k,a}R_{t,k}(\partial_iC_{t,k,b})^\top
  \Big].
  \label{eq:p12-contraction-derivative}
\end{aligned}
\]
This is just the product rule applied to
\(C_{t,k,a}R_{t,k}C_{t,k,b}^\top\), propagated backward from
\(\partial_iR_{t,D}=0\).  Therefore the differentiated contraction is the
derivative of the same direct mass-contraction value \(R_{t,0}\) used in the
declared normalizer.

Core sensitivities come from the fixed core-construction equations.  For an
exact fixed interpolation equation
\[
  A_{t,k}(\alpha)g_{t,k}(\alpha)=b_{t,k}(\alpha),
  \label{eq:p12-fixed-interpolation}
\]
differentiation gives
\[
  A_{t,k}\partial_i g_{t,k}
  =
  \partial_i b_{t,k}
  -
  (\partial_i A_{t,k})g_{t,k}.
  \label{eq:p12-fixed-interpolation-derivative}
\]
because
\[
  \partial_i(A_{t,k}g_{t,k})
  =
  (\partial_iA_{t,k})g_{t,k}+A_{t,k}\partial_i g_{t,k}
  =
  \partial_i b_{t,k}.
\]
For a fixed weighted least-squares update
\[
  g_{t,k}
  =
  \arg\min_g
  \|W_{t,k}^{1/2}(A_{t,k}g-b_{t,k})\|_2^2,
\]
with fixed \(W_{t,k}\) and full-rank
\(N_{t,k}=A_{t,k}^\top W_{t,k}A_{t,k}\), the normal equations are
\[
  A_{t,k}^\top W_{t,k}A_{t,k}g_{t,k}
  =
  A_{t,k}^\top W_{t,k}b_{t,k}.
\]
Differentiating them gives
\[
\begin{aligned}
  N_{t,k}\partial_i g_{t,k}
  =
  &(\partial_i A_{t,k})^\top W_{t,k}
    (b_{t,k}-A_{t,k}g_{t,k})\\
  &+
  A_{t,k}^\top W_{t,k}
  (\partial_i b_{t,k}-(\partial_iA_{t,k})g_{t,k}).
  \label{eq:p12-ls-derivative}
\end{aligned}
\]
The right hand side is computable from the fixed branch once the derivative of
the target square-root evaluations is known.  It is also the derivative of the
same least-squares update, because it is obtained by differentiating the
normal equations that define that update.

The target square-root evaluations depend recursively on the previous
approximate filter.  From
\[
  q_t
  =
  \widehat p_{t-1}^{\TT}
  f_\alpha
  g_\alpha,
\]
where the arguments are those in \eqref{eq:p12-zc-q},
\[
  \partial_i\log q_t
  =
  \partial_i\log\widehat p_{t-1}^{\TT}
  +
  \partial_i\log f_\alpha
  +
  \partial_i\log g_\alpha.
  \label{eq:p12-recursive-logq-derivative}
\]
If
\[
  \widehat p_{t-1}^{\TT}(v;\alpha)
  =
  a_{t-1}(v;\alpha)/\widehat z_{t-1}(\alpha),
\]
then
\[
  \partial_i\log\widehat p_{t-1}^{\TT}(v;\alpha)
  =
  \frac{\partial_i a_{t-1}(v;\alpha)}{a_{t-1}(v;\alpha)}
  -
  \frac{\partial_i\widehat z_{t-1}(\alpha)}
       {\widehat z_{t-1}(\alpha)}.
  \label{eq:p12-prev-filter-derivative}
\]
The numerator \(a_{t-1}\) is a marginal contraction of the previous squared-TT
object.  To see the derivative explicitly, write the previous marginal
numerator as
\[
  a_{t-1}(v;\alpha)
  =
  \int
  \{\phi_{t-1}(v,w;\alpha)^2
  +\tau_{t-1}(\alpha)\lambda_{t-1}(v,w)\}
  \,\dd w,
  \label{eq:p12-prev-marginal-numerator}
\]
where \(w\) denotes the coordinates eliminated from the previous joint
approximation.  Under the same fixed-branch differentiation assumptions,
\[
  \partial_i a_{t-1}(v;\alpha)
  =
  2\int
  \phi_{t-1}(v,w;\alpha)
  \partial_i\phi_{t-1}(v,w;\alpha)
  \,\dd w
  +
  \partial_i\tau_{t-1}(\alpha)
  \int\lambda_{t-1}(v,w)\,\dd w,
  \label{eq:p12-prev-marginal-derivative}
\]
with the corresponding \(\partial_i\lambda_{t-1}\) term if the defensive
density itself is parameter-dependent.  In TT coordinates,
\eqref{eq:p12-prev-marginal-derivative} is evaluated by the same marginal
contraction used for \(a_{t-1}\), except that one core at a time is replaced
by its differentiated core, exactly as in
\eqref{eq:p12-tt-product-rule} and
\eqref{eq:p12-contraction-derivative}.  Thus the derivative is recursive in
time in the same way that the filtering density is recursive.

Combining \eqref{eq:p12-score-general},
\eqref{eq:p12-contraction-derivative}, and
\eqref{eq:p12-prev-filter-derivative}, the direct mass-contraction version of
the same-scalar score is
\[
  \boxed{
  \partial_i\widehat\ell_T^{\TT}
  =
  \sum_{t=1}^T
  \left[
  \frac{\partial_iR_{t,0}+\partial_i\tau_t}
       {R_{t,0}+\tau_t}
  -
  \partial_i c_t
  \right].
  }
  \label{eq:p12-final-score}
\]
This is a full analytical gradient for the declared fixed-branch value path,
provided the selected core-construction branch and all previous-filter
contractions are differentiated by the displayed equations.
\end{proof}

\begin{corollary}[Adaptive code is not globally differentiated by this proof]
If TT-cross pivots, ranks, random/debug samples, truncation decisions,
preconditioner branches, QR/SVD active subspaces, selected minimizers for
\(c_t\), or reapproximation triggers change with \(\alpha\), the formulas above
remain derivatives of the frozen local branch only.  They are not a global
ordinary derivative of the adaptive companion code.
\end{corollary}

\begin{proof}
The proof of Proposition 2 assumes fixed differentiable equations.  A changed
pivot set, rank, sample array, truncation, or selected minimizer changes the
equations defining the scalar.  The derivative of one branch is therefore not
the derivative of a single globally smooth map across branch changes.
\end{proof}

\section{Implementation Checklist}

A future BayesFilter implementation of the fixed-branch variant must specify:
\begin{longtable}{@{}p{0.29\linewidth}p{0.61\linewidth}@{}}
\toprule
Object & Required specification \\
\midrule
\endhead
Filtering model & \(p_0\), \(f_\alpha\), \(g_\alpha\), dimensions of
\(x_t\), \(\vartheta\), and \(u_t=(x_t,\vartheta,x_{t-1})\). \\
Approximation branch & Fixed basis, domains, ordering, ranks, interpolation
or least-squares equations, local solve schedule, and all sample arrays. \\
Approximate density & \(\phi_t^2+\tau_t\lambda_t\), with \(\lambda_t\)
normalized and \(\widehat z_t>0\). \\
Filtering output & Normalized joint \(\widehat\pi_t\) and marginal
\(\widehat p_t^{\TT}(x_t,\vartheta)\). \\
Scalar & Exactly \(\sum_t\{\log\widehat z_t-c_t\}\), or a renamed scalar if
\(c_t\) or the normalizer path changes. \\
Derivative & Core sensitivities, previous-filter sensitivities,
\(\partial_iR_{t,0}\), \(\partial_i\tau_t\), and \(\partial_i c_t\). \\
Diagnostics & Branch stability under finite differences, positive finite
normalizers, rank stability, and same-scalar value-gradient parity. \\
\bottomrule
\end{longtable}

\section{Finite-Difference Parity Required Before HMC Use}

For each parameter coordinate \(\alpha_i\), a future implementation should
check that the value path and derivative path are identical.  With branch
objects frozen,
\[
  \frac{\widehat\ell_T^{\TT}(\alpha+h e_i)
        -\widehat\ell_T^{\TT}(\alpha-h e_i)}{2h}
  \approx
  \partial_i\widehat\ell_T^{\TT}(\alpha).
\]
If the finite-difference calls change rank, pivot, sample, sign, truncation, or
selected constant branches, the check should report branch instability rather
than treating the gradient as an HMC-ready object.

\section{Source And Code Anchors}

The source anchors used here are deliberately narrow.

\begin{itemize}[leftmargin=*]
\item Zhao--Cui equations (1)--(3): state-space transition, observation, and
  parameter notation \cite{zhao-cui-2024}.
\item Zhao--Cui equations (9)--(12) and Algorithm 1: recursive posterior
  update, nonseparable sequential density
  \(q_t(x_t,\theta,x_{t-1})\), TT approximation, and normalizer construction
  \cite{zhao-cui-2024}.
\item Zhao--Cui equation (13) and Lemma 1 vicinity: squared-TT defensive
  approximation \(\widehat\pi=\phi^2+\tau\lambda\) and its normalizer
  \cite{zhao-cui-2024}.
\item Zhao--Cui Proposition 2: marginal density construction from a normalized
  squared-TT approximation \cite{zhao-cui-2024}.
\item Zhao--Cui Algorithm 2: sequential estimation using squared-TT
  approximations \cite{zhao-cui-2024}.
\item Zhao--Cui Section 4.1: exact joint normalizer
  \(z_t=\int\pi_t(x_t,\theta,x_{t-1}\mid y_{1:t})\,dx_t\,d\theta\,dx_{t-1}\)
  is the evidence \(p(y_{1:t})\), and the algorithm approximates \(q_t\)
  because the exact recursion is not directly evaluable \cite{zhao-cui-2024}.
\item Zhao--Cui Proposition 4: conditional KR maps from the squared-TT
  marginal/conditional construction \cite{zhao-cui-2024}.
\item Cui--Dolgov Sections 2.2--2.4 and Section 3: inverse Rosenblatt
  transport, functional TT representation, TT-based inverse Rosenblatt
  construction, and squared inverse Rosenblatt transport substrate used to
  interpret conditional KR maps \cite{cui-dolgov-2022}.
\item Companion code snapshot:
  \texttt{models/full\_sol.m} updates the scalar by
  \texttt{log(sirt.z)-const}, and
  \texttt{deep-tensor.dev/src/@TTSIRT/marginalise.m} sets
  \texttt{obj.z = obj.fun\_z + obj.tau}.
\end{itemize}

\section{What Is Not Proved}

This note does not prove posterior accuracy of the TT approximation, global
differentiability of adaptive TT-cross/rank-changing code, HMC convergence,
production readiness, GPU/XLA readiness, or a default-method recommendation.
It proves two narrower statements: normalized approximate filtering for a
declared fixed branch, and exact same-scalar differentiation for that declared
branch.

\begin{thebibliography}{9}
\bibitem{zhao-cui-2024}
Zhao, Y. and Cui, T. (2024).
\newblock Tensor-train methods for sequential state and parameter learning in
state-space models.
\newblock \emph{Journal of Machine Learning Research}, 25(244):1--51.

\bibitem{cui-dolgov-2022}
Cui, T. and Dolgov, S. (2022).
\newblock Deep composition of tensor-trains using squared inverse Rosenblatt
transports.
\newblock \emph{Foundations of Computational Mathematics}, 22:1863--1922.
\newblock DOI: 10.1007/s10208-021-09537-5.
\end{thebibliography}

\end{document}
